\documentclass[12pt]{article}
\usepackage[utf8]{inputenc}
\usepackage[T1]{fontenc}
\usepackage[margin=1in]{geometry}
\usepackage{setspace}
\usepackage{xcolor}
\usepackage[colorlinks=true, urlcolor=blue!60!black, linkcolor=blue!60!black]{hyperref}
\setlength{\parindent}{2em}
\setlength{\parskip}{1.25em}
\renewcommand{\baselinestretch}{1.0}
\begin{document}
\begin{center}
{\LARGE \textit{Zelman v. Simmons-Harris}, 536 U.S. 639 (2002)\par}
\vspace{0.8em}
{\large Case Brief --- Four-Page Edition\par}
\end{center}
\vspace{1em}

\section{Memory Jogger}

Upholding Cleveland's school-voucher program 5--4, the Court held that a religion-neutral aid program giving assistance to a broad class of families, who direct it to religious schools through their own genuine and independent private choice, does not violate the Establishment Clause (\href{https://supreme.justia.com/cases/federal/us/536/639/}{Justia}).

\section{Detailed Case Facts}

In 1995 a federal court placed the Cleveland City School District, serving over 75,000 mostly low-income and minority students, under state control after it failed all 18 state performance standards. Ohio's Pilot Project Scholarship Program offered tuition scholarships for K--8 students to attend participating private schools, religious or not, or public schools in adjacent districts, and tutoring grants for students who stayed in Cleveland public schools. Low-income families received 90\% of tuition up to \$2,250, with copayments capped at \$250; checks went to parents, who endorsed them to the school. In 1999--2000, 46 of 56 participating private schools (82\%) were religious, no suburban public school participated, and 96\% of more than 3,700 scholarship students attended religious schools. Cleveland also offered nonreligious community (charter) and magnet schools.

\section{Procedural History}

After the Ohio Supreme Court rejected federal Establishment Clause claims but struck the statute on state procedural grounds, which the legislature cured, taxpayers sued in federal court in 1999. The district court granted them summary judgment, and a divided Sixth Circuit affirmed, finding that the program had the primary effect of advancing religion.

\section{Judicial Votes}

Reversed, \textbf{5--4} (decided June 27, 2002).
\begin{itemize}
\item \textbf{Majority:} Chief Justice Rehnquist, joined by Justices O'Connor, Scalia, Kennedy, and Thomas.
\item \textbf{Concurring:} Justices O'Connor and Thomas.
\item \textbf{Dissenting:} Justice Stevens; Justice Souter, joined by Justices Stevens, Ginsburg, and Breyer; Justice Breyer, joined by Justices Stevens and Souter.
\end{itemize}

\section{Holding}

A government aid program that is neutral toward religion and provides assistance directly to a broad class of citizens, who direct it to religious schools wholly through their own genuine and independent private choice, does not violate the Establishment Clause.

\section{Analysis of Opinions}

\textbf{Majority (Rehnquist).} Under \textit{Agostini}, the program's secular purpose was undisputed, so the question was effect. Following \textit{Mueller}, \textit{Witters}, and \textit{Zobrest}, the Court distinguished direct aid from programs of true private choice, in which ``the circuit between government and religion [is] broken'' and any advancement of religion is attributable to parents. The program was neutral, defined its beneficiaries by residence and need, and actually created financial disincentives to choose religious schools, which received less aid than community or magnet schools and required copayments. Genuine choice must be judged across all the options Ohio offered; counting community, magnet, and tutored public students, religious enrollment fell below 20\%. \textit{Nyquist}, involving grants only to private-school parents, did not govern.

\textbf{O'Connor (concurring).} The decision does not break from precedent or from longstanding indirect aid such as Pell Grants and the G.I. Bill.

\textbf{Thomas (concurring).} The Establishment Clause, originally a federalism restraint, may not bind States identically; invoking Frederick Douglass, he called school choice emancipation for urban minority children.

\textbf{Stevens (dissenting).} The educational crisis and parental voluntariness are irrelevant to whether the State may fund religious indoctrination.

\textbf{Souter (dissenting).} The majority abandoned \textit{Everson}'s no-aid principle for ``verbal formalism.'' The tuition cap favored parish-subsidized schools, leaving most recipients a Hobson's choice.

\textbf{Breyer (dissenting).} Funding criteria will entangle officials in religious doctrine and spark sectarian conflict.

\textbf{Later.} \textit{Locke v. Davey} (2004) allowed excluding devotional degrees, but \textit{Trinity Lutheran} (2017), \textit{Espinoza} (2020), and \textit{Carson v. Makin} (2022) turned \textit{Zelman}'s shield into a Free Exercise sword.

\section{Examples: Future Applications}

\begin{enumerate}
\item \textbf{Same-Side (Upheld): Rural Savings Accounts.} Education savings accounts flow entirely to evangelical academies in a county with no secular private schools; local tallies do not matter.
\item \textbf{Same-Side (Exclusion Barred): Excluding Religious Schools.} A State that drops religious schools from that program violates the Free Exercise Clause under \textit{Espinoza} and \textit{Carson}.
\item \textbf{Opposite-Side (Invalid): One Denomination.} Vouchers redeemable only at one church's schools are not neutral.
\item \textbf{Opposite-Side (Invalid): Direct Grants.} State-set payments to religious schools for religion classes involve no private choice.
\item \textbf{Fence-Sitter: Rehabilitation Programs.} Offenders choose among mostly faith-based programs, the secular one waitlisted for six months---arguably a Hobson's choice, with added coercion concerns.
\end{enumerate}

\section{Critique}

\textbf{The denominator problem.} The outcome hinged on the baseline: counting magnet and community schools cut religious participation from 96\% to under 20\%, while Souter's voucher-only denominator showed heavy concentration. Framing the baseline often dictates the result.

\textbf{The vanishing play in the joints.} Together, \textit{Zelman} and \textit{Carson} leave States a binary choice: fund only public schools, or fund secular and religious private schools equally.

\section{Key Quotations}

\begin{quote}
``[W]here a government aid program is neutral with respect to religion, and provides assistance directly to a broad class of citizens who, in turn, direct government aid to religious schools wholly as a result of their own genuine and independent private choice, the program is not readily subject to challenge under the Establishment Clause.''
\end{quote}
\noindent --- Chief Justice Rehnquist, for the Court, 536 U.S. at 652.

\begin{quote}
``The constitutionality of a neutral educational aid program simply does not turn on whether and why, in a particular area, at a particular time, most private schools are run by religious organizations, or most recipients choose to use the aid at a religious school.''
\end{quote}
\noindent --- Chief Justice Rehnquist, for the Court, 536 U.S. at 658.

\begin{quote}
``The criterion is one of genuinely free choice on the part of the private individuals who choose, and a Hobson's choice is not a choice, whatever the reason for being Hobsonian.''
\end{quote}
\noindent --- Justice Souter, dissenting, 536 U.S. at 707.

\end{document}
